\documentclass[11pt,a4paper]{article}

\IfFileExists{glyphtounicode.tex}{\input{glyphtounicode}\pdfgentounicode=1}{}

\usepackage[utf8]{inputenc}
\usepackage[T1]{fontenc}
\usepackage[margin=25mm]{geometry}
\usepackage{amsmath}
\usepackage{amssymb}
\usepackage{booktabs}
\usepackage{longtable}
\usepackage{array}
\usepackage{xcolor}
\usepackage{hyperref}
\IfFileExists{microtype.sty}{\usepackage[protrusion=true,expansion=false]{microtype}}{}

\hypersetup{
  colorlinks=true,
  linkcolor=black,
  citecolor=black,
  urlcolor=blue,
  pdfauthor={Harun Aktas},
  pdftitle={DAD FieldWorks: An Evidence-Controlled Full-Wave CEM Solver Platform Direction},
  pdfsubject={Technical public specification draft},
  pdfkeywords={computational electromagnetics, evidence contracts, full-wave CEM, residuals, eigenpairs, claim boundaries}
}

\setlength{\parindent}{0pt}
\setlength{\parskip}{0.62em}
\renewcommand{\arraystretch}{1.18}
\emergencystretch=2em

\newcommand{\code}[1]{\texttt{#1}}
\newcommand{\statusbox}[1]{%
  \vspace{0.55em}
  \noindent\fcolorbox{black}{gray!7}{%
    \begin{minipage}{0.94\linewidth}
    #1
    \end{minipage}}
  \vspace{0.55em}
}

\title{\textbf{DAD FieldWorks}\\
{\large An Evidence-Controlled Full-Wave CEM Solver Platform Direction}\\
\vspace{0.4em}
{\normalsize Evidence Contracts for Candidate Fields, Eigenpairs, Residuals and Claim-Bounded Electromagnetic Results}}
\author{Harun Aktas}
\date{Technical Whitepaper / Public Specification Draft\\Version 0.9 draft\\June 2026}

\begin{document}

\maketitle

\statusbox{\textbf{Document status:} Public architecture draft.\\
\textbf{Claim status:} No external validation claim. No production readiness claim. No commercial solver equivalence claim.\\
\textbf{Copyright:} Copyright \copyright{} 2026 Harun Aktas. All rights reserved.\\
\textbf{License notice:} No license is granted for reuse, redistribution, or derivative work without written permission from the copyright holder.}

\begin{abstract}
DAD FieldWorks is being developed toward an evidence-controlled full-wave computational electromagnetics solver platform. The project separates numerical output from what that output is allowed to claim. Current public-safe material remains bounded internal alpha evidence and diagnostic visualization only. Residual, analytical comparison and physical eigenpair acceptance gates are treated as evidence layers, not production validation. No external validation or production readiness is claimed.
\end{abstract}

\textbf{Keywords:} computational electromagnetics; evidence contracts; full-wave CEM; eigenpairs; residual checks; reference comparison; claim boundaries.

\tableofcontents
\newpage

\section{Introduction}

Field plots, eigenvalues and spectra can look plausible while still being wrong. A trustworthy computational electromagnetics workflow must expose residuals, boundary assumptions, reference comparisons, reproducibility metadata and claim boundaries. DAD FieldWorks is being developed around this discipline: a field, eigenpair or spectrum remains a candidate claim until the evidence record defines what it is allowed to say.

This document is a public specification draft for a private research and development project. It presents a public-safe architecture direction and current bounded alpha evidence layers. It does not publish private solver source code, private workbench material, internal project-control files or private validation artifacts.

\section{Why Field Plots Are Not Enough}

A computed field is not trustworthy because it is visually convincing. A candidate eigenpair is not accepted merely because it exists. A spectral peak is not a validation result merely because it is visible.

For full-wave computational electromagnetics, the missing evidence may include residual quality, component-location policy, boundary compatibility, grid convention, reference-problem equivalence, reproducibility records and explicit claim authorization. DAD FieldWorks therefore treats visual and numerical outputs as candidate records until evidence gates support a bounded interpretation.

\section{Evidence-Controlled Full-Wave CEM Solver Direction}

The long-term goal is a native full-wave solver platform where solver outputs carry evidence records by default. A field, eigenpair or spectrum is not promoted because it looks plausible. It remains claim-bounded until residual, reference, reproducibility and validation evidence support a stronger statement.

This is a direction, not a completed product statement. The current public material does not state that DAD FieldWorks is already a full-wave production solver.

\section{Evidence Contract Model}

An evidence contract couples a numerical payload with evidential state. The payload may be a field, candidate eigenpair, spectrum, trace, frequency estimate, residual, diagnostic row or replay artifact. The contract records the assumptions, checks, gates, blockers and claim boundaries that govern interpretation.

\begin{description}
  \item[Candidate Result] A computed or replayed result that is not yet promoted beyond its evidence record.
  \item[Evidence Row] A factual row about geometry, grid, field component, residual value, finite-value status, comparison input or replay provenance.
  \item[Gate Row] A pass, fail, blocked or not-applicable row that controls whether a candidate can support a stronger statement.
  \item[Residual Row] A row recording residual definition, residual magnitude, normalization and classification scope.
  \item[Reference Comparison Row] A row binding an internal result to a stated reference problem, reference value, comparison metric and equivalence policy.
  \item[Acceptance Gate] A gate that may accept, reject or block a candidate only under explicit prerequisites.
  \item[Blocker Row] A row preserving why promotion is not allowed, such as missing prerequisites or unresolved mapping policy.
  \item[Claim Boundary] A statement of what must not be inferred from the current evidence.
  \item[\code{ProductionAllowedQ}] A conservative flag that remains false unless a separate production gate supports changing it.
  \item[\code{ExternalValidationQ}] A conservative flag that remains false unless real external validation evidence and gates support changing it.
\end{description}

Numerical success does not automatically become a validation claim.

\section{Current Public-Safe Alpha State}

The internal alpha path has progressed through bounded C++ Resonator Lab Alpha record layers for generalized eigenproblem metadata, candidate eigenpair records, residual diagnostics, residual threshold gate records and analytical comparison gate records. The next current direction is records-only physical eigenpair acceptance gate planning.

This state is public-safe only at the summary level. It should be read as bounded internal alpha evidence, not as a release of implementation details and not as public validation evidence.

\section{C++ Resonator Lab Alpha Evidence Path}

The C++ Resonator Lab Alpha direction is record-level and evidence-contract oriented. Publicly described elements include generalized eigenproblem metadata, candidate eigenpair record handling, residual diagnostics, JSON and CSV style evidence artifacts and bounded internal alpha status.

This does not claim a production C++ electromagnetic solver, production numerical kernel, production eigensolver or accepted physical eigenfrequency. It is a disciplined record path for candidate results and replayable evidence.

\section{Candidate Eigenpair and Residual Diagnostics}

A candidate eigenpair is treated as a claim. Residual calculation provides evidence about whether the candidate is compatible with the stated operator and mass model under the current assumptions. Residual magnitude can support classification, but it does not automatically become physical acceptance.

The public evidence-gate visualizations use this principle: a candidate is computed or displayed, the residual is reported, and the result remains inside a claim boundary unless the required gates justify promotion.

\section{Residual Threshold and Analytical Comparison Gates}

Residual threshold gate logic and analytical comparison records are evidence gates. Their purpose is to classify or block candidate records under explicit prerequisites.

The residual threshold gate asks whether the residual evidence is strong enough under the current normalization and prerequisite policy. The analytical comparison gate asks whether a candidate result is close enough to a stated analytical reference under an explicit comparison policy. These gates are not external validation.

\section{Physical Eigenpair Acceptance Gate Planning}

The current next direction is planning a physical eigenpair acceptance gate. This means planning schemas, prerequisite evidence requirements, residual threshold pass requirements, analytical comparison pass requirements, acceptance metadata, pass/fail and blocker records, and future audit boundaries.

The physical eigenpair acceptance gate is planning-only in the public claim boundary. No physical eigenpair has been accepted by public claim. No physical eigenfrequency has been accepted by public claim.

\section{Public Showcase Diagnostics}

The current featured public visual diagnostics are:

\begin{itemize}
  \item PEC Resonator Candidate Evidence Gate,
  \item Eigenpair Residual Evidence Gate.
\end{itemize}

They are public communication assets for candidate, residual, evidence-gate and claim-boundary concepts. \code{ProductionAllowedQ} remains false. \code{ExternalValidationQ} remains false. These diagnostics do not claim physical eigenpair acceptance, physical eigenfrequency acceptance, external validation or production readiness.

\section{Long-Term Solver Roadmap}

\begin{center}
\begin{tabular}{p{0.27\linewidth}p{0.34\linewidth}p{0.29\linewidth}}
\toprule
Stage & Direction & Public Claim Boundary \\
\midrule
Current alpha evidence layer & Candidate, residual, analytical comparison records & Bounded internal alpha evidence only \\
Current planning step & Physical eigenpair acceptance gate planning & Planning only, no physical acceptance execution \\
Near term & Residual threshold and acceptance gate hardening & Internal classification only \\
Mid term & Native solver core hardening and artifact replay & No production solver claim \\
Long term & Evidence-controlled full-wave CEM solver platform & Future target, no current production claim \\
Future research route & DGTD and physics-AI backends under evidence contracts & Future route only \\
\bottomrule
\end{tabular}
\end{center}

The roadmap is evidence-led rather than feature-led. A future backend is not trusted because it runs; it is trusted only within the evidence boundary that its gates support.

\section{Limitations and Claim Boundaries}

The public claim boundary is deliberately narrow:

\begin{itemize}
  \item no external validation claim,
  \item no production readiness claim,
  \item no commercial solver equivalence claim,
  \item no production full-wave solver claim,
  \item no physical eigenpair acceptance claim,
  \item no physical eigenfrequency acceptance claim,
  \item no mode identification claim,
  \item no rank/nullspace claim,
  \item no CPML support claim,
  \item no graphical user interface claim,
  \item no qubit simulation claim,
  \item no complete quantum hardware solver claim,
  \item no measurement validation claim.
\end{itemize}

These limitations are part of the architecture. A claim remains closed until the corresponding evidence path exists and passes.

\section{Conclusion}

DAD FieldWorks is being developed toward an evidence-controlled full-wave CEM solver platform. Its purpose is not merely to produce field plots, spectra or eigenpair candidates, but to bind such artifacts to residuals, reference comparisons, reproducibility records, gate decisions, failure classification and claim boundaries.

The current public state is useful but intentionally bounded: public-safe alpha evidence paths and diagnostic visualizations only. No external validation claim is made. No production readiness claim is made. No commercial solver equivalence claim is made.

\section{References}

\begin{thebibliography}{9}

\bibitem{yee1966}
Kane S. Yee,
``Numerical solution of initial boundary value problems involving Maxwell's equations in isotropic media,''
\emph{IEEE Transactions on Antennas and Propagation}, 1966.

\bibitem{roache}
Patrick J. Roache,
\emph{Verification and Validation in Computational Science and Engineering},
Hermosa Publishers, 1998.

\bibitem{oberkampf}
William L. Oberkampf and Christopher J. Roy,
\emph{Verification and Validation in Scientific Computing},
Cambridge University Press, 2010.

\bibitem{taflove}
Allen Taflove and Susan C. Hagness,
\emph{Computational Electrodynamics: The Finite-Difference Time-Domain Method},
Artech House, third edition, 2005.

\end{thebibliography}

\end{document}
